\section{Introduction and the state of the problem}
\label{sec:introduction}

The ProveIt polylogarithm manuscript brings together exact cyclotomic
identities, positive-kernel representations, and zero geometry of
Stieltjes--Lerch derivatives. The most useful next step is to distinguish
questions about a fixed integer weight from questions about an entire
family. A special-value reduction asks which periods suffice to express
one number. An order interpolation asks how those numbers fit into a
holomorphic function. A zero-motion problem asks what remains uniform
when two discrete indices change together. These questions interact, but
a solution to one does not automatically solve the others.

This article develops three results beyond the inspected source baseline.
The proofs are given in full, with classical analytic inputs either
derived here or explicitly identified. The accompanying programs check
finite identities and numerical conventions; the infinite conclusions
do not rest on an integer-relation search or a scan of zeros.

\subsection{The harmonic Euler sum as a function of its order}

Let $H_0=0$ and $H_n=\sum_{j=1}^n j^{-1}$. The manuscript's mixed
Gaussian sums are the positive-integer values of
\[
 S(s)=\sum_{n=0}^{\infty}\frac{(-1)^nH_n}{(2n+1)^s},
 \qquad \Re s>0.
\]
Our first main result is that this function extends to an entire function
whose complete real zero set is
\[
 \{\sigma_m:m\ge1\},\qquad -2m-1<\sigma_m<-2m,
\]
with every zero simple. In particular $S(s)<0$ for real $s\ge-2$.
We also prove that $S$ has no complex zero in $\Re s\ge3$.

The continuation itself is a standard Mellin-regularization argument.
The new work is the global zero classification. A differentiated
Fourier--beta integral gives
\[
 S(-q)=\frac{cD(q)}2\{R(q)\cos(cq)-\sin(cq)\},\qquad c=\frac\pi2,
\]
where $D(q)>0$ and $R$ is a ratio of explicit Mellin moments.
We prove $R(q)>0$ and $0<R'(q)<\pi/2$ for $q\ge2$. This converts
existence and uniqueness in every interval into a strict comparison
with the tangent function. The derivative estimate follows from a
one-dimensional variance inequality and elementary moment bounds.
It covers the first interval as well as the asymptotic regime.

If $N=2m+1$, $\Lambda=\log(N/\pi)+\gamma$, and
$\sigma_m=-N+\varepsilon_m$, then
\[
 \varepsilon_m=
 \frac2\pi\arctan\frac{\pi}{2\Lambda}
 -\frac{\frac12-\frac2\pi\arctan(\pi/(2\Lambda))}
        {N(\Lambda^2+\pi^2/4)}
 +O\!\left(\frac1{N^2\Lambda^2}\right).
\]
We give both a complete expansion in inverse logarithms and a recursive
construction of every correction in $1/N$. Exact negative-integer
identities in generalized secant numbers supply an independent arithmetic
description.

\subsection{Diagonal Lerch zeros do not have an eventual common direction}

For $n\ge0$, $k\ge1$, $h>0$, and $0\le\rho\le1$, define
\[
 \mathscr F_{n,k}^{\rho,h}(a)=
 \sum_{m=0}^{\infty}\rho^m(-1)^{n+k}
 \frac{d^k}{da^k}\left(\frac{\log^n(a+mh)}{a+mh}\right),
 \qquad a>0.
\]
The endpoint $\rho=1$ is defined by this convergent derivative series.
On the diagonal $k=n-1$, the smallest elementary zero is $a=1$.
It continues to an analytic local zero branch $a_n^{(h)}(\rho)$.
For $A=1+h$, put $v_n(A)=(a_n^{(h)})'(0)$.
We prove the exact convergent generating identity
\[
 \sum_{n\ge1}v_n(A)t^n=\log A-u_A(t),\qquad
 e^{u_A(t)}=A+t\,u_A(t),\qquad u_A(0)=\log A.
\]
Here the $n=1$ coefficient is an auxiliary term completing the generating
function; the zero-branch assertion concerns $n\ge2$.

For every real $A>1$, infinitely many $v_n(A)$ are positive and infinitely
many are negative. The proof uses Lagrange inversion and a
positive-coefficient argument: the inverse $u_A$ is analytic along the
whole positive real axis and has unbounded logarithmic growth, which
excludes eventually one-sign Taylor coefficients. This is compatible
with a theorem that fixes $n$ before letting $k$ grow; it shows why that
order of quantifiers matters.

\subsection{An explicit solution to the exterior-operator problem}

For $n\ge2$, set
\[
 g_n(x)=x^{-2}(\log x)^{n-1}(n-\log x),\qquad
 U_n(a,\rho)=\sum_{m\ge0}\rho^m g_n(a+m).
\]
The incoming \emph{Lerch Global Phase} report asks for the best exterior
cutoff obtainable by optimizing the positive constant in
$\partial_a+\lambda$ \cite{IncomingPhase}. Let $r_n$ be the unique
positive root of
\[
 r^3+nr^2-nr-n=0.
\]
The optimal coefficient is
\[
 \lambda_n=(2-r_n)e^{-n-r_n},\qquad 1<r_n<\frac{1+\sqrt5}{2}.
\]
We give a scalar equation uniquely defining the least cutoff $B_n$,
prove
\[
 (\partial_a+\lambda_n)U_n(a,\rho)<0
 \quad(a\ge B_n,\ 0<\rho\le1),
\]
and prove that no smaller cutoff works uniformly over these parameters.
At that cutoff the coefficient is unique. For the corresponding strict
pointwise problem, the same cutoff is an infimum and is not attained:
there is an unavoidable tangency in the negative tail. This boundary
qualification is part of the solution.

The algebraic parameter $r_n$ tends to the golden ratio. Its series in $1/n$ has exact
radius of convergence $5/27$, and the cutoff has explicit exponentially
small corrections in logarithmic coordinates. These are quantitative
identities, rather than numerical fits.

\subsection{Source baseline, attribution, and proof status}
\label{sec:baseline}

The source baseline is the repository commit
\[
 \texttt{a0a90ef31877f98be437191c48b46f02d5456867},
\]
dated 10 October 2026. The focused audit used all forty manuscript chapter
files and the mathematical source material in the sixteen incoming
archives listed in the accompanying provenance record. This is a focused
research audit, not a claim to have independently re-proved every theorem
in those archives.

At this baseline, the mixed Gaussian $S_4$ reduction is already proved
by a convergent octahedral transformation and an exact finite certificate
\cite{ProjectS4}. The $S_2$ companion, the product-matrix rank theorem,
and several low-index Lerch phase results are also already established.
The two displayed $S_6$ baskets are provably equivalent formulations of
one remaining conjecture \cite{ProjectQuotients}. A frozen $S_8$ vector
has already been disproved by an exact interval exclusion
\cite{ProjectDiscovery}. None of those earlier achievements is claimed
as a result of this article.

Analytic continuation of harmonic Euler Dirichlet series has substantial
precedent \cite{BGP}; integer-weight Euler sums and their colored
multiple-zeta representations are also established subjects
\cite{XuWang}. Lagrange inversion, the Pringsheim principle, and the
Brascamp--Lieb inequality are classical \cite{Gessel2016,Flajolet,BL}.
The claims of progress here are relative to the inspected manuscript and
incoming reports. The literature check was targeted and does not establish
global priority for every formula.

The results designated as theorems below have ordinary analytic or
algebraic proofs. Exact rational computations provide additional finite
certificates. Numerical quadrature and high-precision asymptotic
comparisons are labelled diagnostics. The complex-singularity conjecture
is kept separate from the proved infinite sign theorem, and $S_6$ remains
open.

\begin{table}[htbp]
\centering
\small
\begin{tabularx}{\textwidth}{@{}p{.28\textwidth}X p{.19\textwidth}@{}}
\toprule
Question & Result in this article & Status\\
\midrule
Real zeros of $S(s)$ &
Entire continuation; exactly one simple zero in each
$(-2m-1,-2m)$; no other real zeros & Proved\\
Diagonal initial motion &
Exact inverse generating function; infinitely many velocities of each sign
for every $h>0$ & Proved\\
Best constant exterior operator &
Unique optimal coefficient and least uniform summed cutoff;
strict pointwise infimum identified & Proved\\
Quantitative diagonal oscillation &
Candidate accessible conjugate singularities and coefficient asymptotics
& Explicit conjecture\\
Mixed reduction for $S_6$ &
Existing candidate and its two-basket equivalence retained
& Still open\\
\bottomrule
\end{tabularx}
\caption{The principal proof-status distinctions.}
\label{tab:status}
\end{table}
